\section{实验设计}

\begin{frame}{数据与样本设计}
  \begin{columns}[T,onlytextwidth]
    \column{0.54\textwidth}
    \begin{block}{数据设置}
      \begin{itemize}
        \item 研究对象：A股上市公司
        \item 时间范围：\hlb{2007--2022 年}
        \item 标签定义：下一年是否被 ST
        \item 特征维度：10 个标准化财务变量
      \end{itemize}
    \end{block}

    \begin{exampleblock}{样本规模}
      \begin{tabular}{@{}lrr@{}}
        \toprule
        数据集 & 样本数 & 说明 \\
        \midrule
        训练集 & 17,437 & 模型拟合 \\
        验证集 & 7,568 & 调参与早停 \\
        测试集 & 6,175 & 一次性样本外评估 \\
        \bottomrule
      \end{tabular}
    \end{exampleblock}

    \column{0.42\textwidth}
    \begin{alertblock}{数据特点}
      \begin{itemize}
        \item 类别极度不平衡
        \item 测试集中 ST 占比仅约 \hl{0.63\%}
        \item 更接近真实预警场景，也更考验模型质量
      \end{itemize}
    \end{alertblock}
  \end{columns}
\end{frame}

\begin{frame}{实验流程与评价指标}
  \begin{columns}[T,onlytextwidth]
    \column{0.50\textwidth}
    \begin{block}{实验流程}
      \begin{itemize}
        \item 用 $t-2,t-1,t$ 三期序列预测 $t+1$ 年状态
        \item 采用 \hlb{时间顺序严格分割}，避免未来信息泄露
        \item LSTM 以验证集 AUC 进行早停，最佳轮次回滚
      \end{itemize}
    \end{block}

    \column{0.46\textwidth}
    \begin{exampleblock}{评价指标}
      \begin{itemize}
        \item \metric{AUC}{区分能力}
        \item \metric{KS}{排序与分离度}
        \item \metric{Precision}{预警可靠性}
        \item \metric{Recall}{风险覆盖率}
        \item \metric{F1}{综合平衡}
      \end{itemize}
    \end{exampleblock}
  \end{columns}

  \begin{alertblock}{为什么不用准确率}
    在极度不平衡数据中，简单准确率会掩盖少数类识别能力，因此本文以 \hl{AUC 和 F1} 为核心指标。
  \end{alertblock}
\end{frame}
